\chapter{关于本书}

本书是一份开放教育资源。它面向稍有逻辑学背景和高中数学知识的学生。其目标是浅尝集合论哲学的皮毛。读完本书，学生可能会对以下几点有所认识：

\begin{enumerate}
	\item 集合论是如何产生的；
	\item 如何在集合论加算术的框架下嵌入数学的大部分领域；
	\item 如何将算术本身嵌入集合论；
	\item 集合的累积迭代概念意味着什么；
	\item 如何尝试为 ZFC 的公理辩护。
\end{enumerate}

本书脱胎于我在剑桥大学哲学系讲授的一门短期课程。在我之前，这门课由 Luca Incurvati（卢卡·因库尔瓦蒂）和 Michael Potter（迈克尔·波特）主讲。在撰写本书——以及更广泛地，准备这门课程——的过程中，我深感受惠于 Luca 和 Michael 两人。我希望这份感激之情能在本书的正文中体现出来；但在此，我也要向他们表达由衷的感谢。

本书的大部分内容最初以《Open Set Theory》之名发布；本书是其续作。我已将本书所有材料贡献给 \href{http://openlogicproject.org}{\emph{Open Logic Project}}，在那里可以免费获取。（\Crefrange{sfr:set::chap}{sfr:siz::chap} 取材自 OLP 中已有的材料，仅作了细微改动；这些改动现在也已成为 OLP 的一部分。）欲了解更多信息，请访问 \href{http://openlogicproject.org/}{openlogicproject.org}。

\begin{flushright}
Tim Button（蒂姆·巴顿）
\\伦敦大学学院
\\2021年10月
\end{flushright}